\begin{center}
    \textbf{Experience}
\end{center}
\textbf{ScalAR Lab, University of Pennsylvania} \hfill Philadelphia, PA

\textbf{Graduate Research Assistant} \hfill August 2025 -- Present (On Medical Leave)\\
PI: Ani Hsieh
\begin{itemize}[noitemsep, topsep=0pt, partopsep=0pt, parsep=0pt]
    \item Researched data-driven modeling of the hydrodynamics of underwater robotic arms, toward model-based controllers that anticipate fluid forces rather than relying on stiff, power-hungry high-gain control.

    \item Built a knowledge-based neural ODE (KNODE) model of a 7-joint Franka Emika Panda arm that combines differentiable rigid-body dynamics from MuJoCo MJX with a neural network that learns the residual fluid torques, trained end to end on GPU in JAX, Equinox, and Diffrax.

    \item Compared KNODE with Fourier-basis SINDy, Gaussian process, and nominal physics models on simulated underwater trajectories, where KNODE most closely matched the ground truth.
\end{itemize}

\vspace{12pt}

\textbf{Grauer Lab} \hfill State College, PA

\textbf{Undergraduate Research Assistant} \hfill January 2024 – May 2025\\
PI: Samuel Grauer
\begin{itemize}[noitemsep, topsep=0pt, partopsep=0pt, parsep=0pt]
    \item Applied scientific machine learning to solve inverse problems of flow state estimation based on optical sensing techniques.
    
    \item Developed a neural optical flow algorithm to reconstruct turbulent flow states from digital images of advected particles in fluids.

    \item Embedded a divergence free condition as a hard constraint within the optical flow problem that utilizes automatic differentiation within PyTorch.

    \item Implemented a soft constraint on the Navier-Stokes residuals to infer latent fields such as pressure.

    \item Impacts of this work include the design of aerofoils, medical devices, and heat exchanging technology via superior field reconstruction abilities and latent field inference.
    
    
\end{itemize}

\vspace{12pt}

\textbf{The SHAPE Lab} \hfill State College, PA

\textbf{Undergraduate Research Assistant} \hfill September 2023 – May 2024\\
PI: Guha Manogharan
\begin{itemize}[noitemsep, topsep=0pt, partopsep=0pt, parsep=0pt]

    \item Investigated the properties of material interface within Multi-Material Laser Powder Bed Fusion additive manufacturing and characterized the effect of interlocking material regions on part strength.
    

    
\end{itemize}

\vspace{12pt}

\textbf{Thrive Scholars} \hfill Amherst, MA

\textbf{Summer Academy Teaching Assistant} \hfill June 2022 – August 2022, June 2023 -- August 2023
\begin{itemize}[noitemsep, topsep=0pt, partopsep=0pt, parsep=0pt]
    \item Assisted the teaching of Calculus 1 and Algebra to high-achieving students of color from economically under-resourced communities during the Summer Academy.
    %\item Graded 20+ papers multiple times per week.
    %\item Tutored groups and individual students weekly.
    %\item Held two hours of office hours each week.
    \item Planned and taught lectures on limits, derivatives, parametric equations, and other fundamental mathematical concepts.
\end{itemize}
